% -*- coding: utf-8 -*-
\chapter{总结与展望}
\label{chap:conclusion}

\section{全篇工作总结}

可重构智能表面（RIS）作为 6G 移动通信物理层的主动电磁调控使能技术，为突破毫米波与太赫兹高频通信的严重路径衰减与视距阻挡瓶颈提供了极具前景的解决方案。本论文围绕 RIS 辅助的多天线基站下行多用户通信系统，系统性开展了主被动协同波束赋形优化与低开销级联信道估计关键技术研究。全论文的主要研究工作与学术贡献总结如下：

\begin{enumerate}
  \item \textbf{建立了自洽严谨的 BS-RIS-User 多用户 MISO 系统级形式化理论模型}：
  基于电磁超材料物理特性与空间几何信道模型，形式化推导了直连链路与无源级联反射链路叠加后的复合传输方程。表征了多用户在接收端的瞬时信干噪比（SINR）与全网可达总和速率，为算法设计确立了坚实的数学物理根基。

  \item \textbf{提出了基于二次变换分式规划与流形优化的主被动波束协同解耦优化机制}：
  针对和速率极大化目标函数中发射预编码与无源连续相移的高度非凸非线性耦合难题，利用拉格朗日多比率二次变换将其解耦为两个凸/拟凸可解子问题。推导了基站发射预编码向量的拉格朗日闭式最优解；引入复圆黎曼流形共轭梯度法（RCG）直接在流形切空间求解相移向量，彻底克服了传统半正定松弛（SDR）算法矩阵升维导致的计算复杂度爆炸缺陷。在非凸优化理论框架下严格证明了算法的目标函数值单调收敛性。

  \item \textbf{设计了基于正交 DFT 码本的极低开销级联信道估计协议}：
  针对无源超表面缺乏射频链导致信道无法直接测量的核心瓶颈，设计了基于时域正交导频与截断 DFT 反射矩阵的估计协议。将导频开销压缩至理论极限的 $N+1$ 个子时隙，利用最小二乘（LS）解析估计器直接实现直连链路与级联信道的高信噪比解耦估计，完全避免了矩阵求逆运算。理论推导证明该设计严格达到了 Fisher 信息的克拉美罗下界（CRLB）。

  \item \textbf{构建了符合国标与学术诚信全流程 Definition of Done (DoD) 的脚手架工程体系}：
  在第四章建立了完备的系统级仿真基准参数集与预置标准三线表，明确保持全量 \TODO 占位规范，杜绝任何学术编造，实现了从文献、模型、理论到工程实现的严密闭环。
\end{enumerate}

\section{未来研究工作展望}

尽管本论文在 RIS 辅助 6G 通信的波束优化与信道估计方面取得了一定的理论进展，但面向 6G 实际网络规模商用部署，仍有诸多前沿方向亟待进一步深入探索：

\begin{enumerate}
  \item \textbf{有限精度离散相移与非理想电磁幅相耦合补偿}：
  实际商用超材料微元多采用 1-bit（0/$\pi$）或 2-bit 离散控制，且相移调整常伴随幅度损耗。未来工作可进一步探索量化离散约束下的凸松弛投影优化，以及基于真实散射参数阻抗网络的自适应鲁棒补偿算法。

  \item \textbf{超大规模阵列（XL-RIS）近场球面波传播模型与宽带波束倾斜效应}：
  当超表面物理孔径达到米级量级且载波频率提升至太赫兹时，瑞利距离大幅扩展，通信用户将落入电磁近场（Near-Field）辐射区域。传统平面波远场假设失效，必须引入基于球面波前的高维近场波束聚焦模型，同时攻克宽带正交频分复用（OFDM）系统中的空间波束倾斜（Spatial Beam Squint）问题。

  \item \textbf{基于 RIS 的通信感知一体化 (ISAC) 理论拓展}：
  6G 网络将深度融合通信与高精度雷达感知能力。利用 RIS 不仅重塑通信链路，更可协同形成面向非视距目标的感知波束，实现目标定位与成像。如何在发射功率预算下权衡通信和速率与雷达感知克拉美罗界的帕累托前沿，具有极其重要的学术研究与产业应用价值。
\end{enumerate}
